\chapter{问题、状态与闭环}

\section{信息赤字}

蒸制失误、烤制不均、时间无法事先给定，以及干烧、超温、漏电，都发生在同一个缺口上：腔体能够改变边界，却没有食品内部状态的充分统计量。把缺口写成传感器清单，得到的是功能表；写成信息量，得到的是本文的设计准则。

\begin{definition}[热加工信息赤字]
真实状态 $\vect{x}(t)\in\cX$，可安装传感器给出 $\vect{y}(t)=h(\vect{x}(t))+\vect{v}(t)$。对决策泛函 $\phi$，
\begin{equation}
\Delta\mathcal{I}(t)=H\big(\phi(\vect{x}(t))\mid \vect{y}_{0:t}\big)-H_{\mathrm{tol}}.
\end{equation}
$\phi$ 只包含与停机、表面品质和安全直接相关的量。$H_{\mathrm{tol}}$ 是该量的误判上限。$\Delta\mathcal{I}>0$ 时仍按开环时间表改变边界，冷点进度与表面干耗会分开漂移。
\end{definition}

系统名称 ACTS-MCDT：Autonomous Cognitive Thermal System with Multi-scale Causal Digital Twin。它要完成的链条是

\[
\begin{aligned}
&\text{边界激励}
\to\text{多尺度状态}
\to\text{后验}\\
&\to\text{信息价值与品质价值}
\to\text{安全投影}
\to\text{新的边界}.
\end{aligned}
\]

\section{五个必须同时成立的性质}

\begin{enumerate}
\item \textbf{可重构。} 冷点温度、表面含水、反应进度和内部热阻，能由有限观测与短时边界激励稳定更新。
\item \textbf{可分离。} 同一次阀门动作，要么为了降低后验熵，要么为了推动品质泛函，二者的目标函数不同，触发条件可检查。
\item \textbf{可迁移。} 宏观传热、孔隙水分、蛋白与淀粉反应耦合成低维决策状态；食材类别只改变先验，不改变状态定义。
\item \textbf{可更新。} 单次过程更新传递参数；结构方程的父节点不随数据改写。
\item \textbf{可拦截。} 品质模型的输出在进入执行器之前，被只依赖硬测量的不等式投影；硬保护不经过该模型。
\end{enumerate}

\section{决策量}

用户侧需要四个判断：能否停止，表面是否已偏离目标含水，温度场是否还在散开，继续施加边界是否安全。

\begin{table}[H]
\centering
\caption{四个判断及所需隐藏量}
\begin{tabularx}{\textwidth}{lXX}
\toprule
判断 & 隐藏量 & 信息来源 \\
\midrule
能否停止 & 冷点进度 $\alpha_m$ 的下分位 & 逆估计与探针确认 \\
表面是否偏离 & $M_s$、汁液损失 $L$ & 质量、湿度、冷凝 \\
是否仍不均匀 & 温度离散度 $U_T$ & 热像与降阶场 \\
继续是否安全 & 干烧、漏电、腔压、门封 & 不进入品质模型的硬测量 \\
\bottomrule
\end{tabularx}
\end{table}

腔温只描述空气。新增的每一路观测必须降低上表中至少一个判断的 $\Delta\mathcal{I}$，否则不进入观测向量。

\section{闭环的完整性条件}

一次运行在下列六步闭合。缺任何一步，该运行不进入结果比较。

\begin{enumerate}
\item 真值模型生成食品与腔体边界，参数对在线模块不可见。
\item 在线模块只接收可安装通道的带噪采样。
\item 后验给出 $\vect{z}$、信息赤字和建议执行量。
\item 执行量经安全投影写回真值模型的边界。
\item 评分使用真值场上的进度、均匀性、能耗和越界次数。
\item 探针前后的协方差若未下降，该次动作记为无效信息，不得改变停止时刻。
\end{enumerate}

真值写入在线模块会使第 2 步失效，该次数据作废。

\section{文档结构}

后文按闭环展开，而不是按装置清单展开。

\begin{center}
\begin{tabularx}{\textwidth}{llX}
\toprule
章 & 对象 & 在闭环中的位置 \\
\midrule
2 & 已有模型与产品能力 & 划定本文新增的环节 \\
3 & 层、时间尺度、模块合同 & 规定信息不能走捷径 \\
4 & 多尺度场与决策状态 & 生成被估计的对象 \\
5 & 观测、病态性、滤波 & 被动信息的上限 \\
6 & 对偶控制与探针 & 主动信息 \\
7 & 结构因果与路径 & 解释与干预 \\
8 & 算子与参数更新 & 跨时间的模型 \\
9 & 品质优化与安全投影 & 执行 \\
10 & 真值隔离的试验 & 闭环是否成立 \\
11 & 可实施特征与实施例 & 把公式收成装置上的步骤 \\
12 & 命题、证明与图件 & 给出可否定的陈述和图的坐标 \\
13 & 交互演示与数据发布 & 规定数值核、界面与下载字段 \\
14 & 文献分层 & 每条方程所继承与所不继承的结果 \\
15 & 图件总表 & 机制图与数据图的坐标和归属 \\
16 & 成稿 & 同一正文的 PDF 与 Word \\
17 & 外部数据 & 哪些参数被公开数据约束 \\
18 & 构建顺序 & 从大纲到文档、数据、仿真与演示 \\
\bottomrule
\end{tabularx}
\end{center}
